
\chapter*{Résumé}
\addcontentsline{toc}{chapter}{Résumé}
Les progrès récents de la traduction automatique neuronale ont permis d'améliorer considérablement la qualité des systèmes de traduction pour les langues disposant d'importantes ressources linguistiques. En revanche, les langues à faibles ressources, telles que le ruund, restent peu prises en compte en raison du manque de corpus parallèles et de ressources numériques. Ce mémoire s'inscrit dans cette problématique et vise à étudier les approches de traduction automatique adaptées aux langues à faibles ressources à travers le développement d'un système de traduction neuronale pour le couple ruund–français.

Pour répondre à ce défi, un corpus parallèle, \textbf{OpenRuund}, a été constitué à partir de textes du Nouveau Testament et de textes poétiques, pour un total de 8~431 paires de phrases. Après numérisation, reconnaissance optique de caractères, nettoyage, normalisation et alignement, trois modèles de traduction neuronale, mBART-50, MarianMT et NLLB-200-distilled-600M, ont été adaptés par \textit{fine-tuning}. Les performances ont été évaluées à l'aide des métriques BLEU et chrF, puis complétées par une évaluation humaine menée par des locuteurs natifs du ruund et un spécialiste en linguistique africaine.

Les résultats montrent que les modèles préentraînés peuvent produire des traductions pertinentes malgré le faible volume de données disponibles. Le modèle NLLB-200-distilled-600M obtient les meilleurs résultats dans les deux sens de traduction, avec un score BLEU de 38.82 et un score chrF de 58.10 pour la traduction français$\rightarrow$ruund, ainsi qu'un score BLEU de 24.11 et un score chrF de 54.80 pour la traduction ruund$\rightarrow$français. Il a été intégré dans une application web associée à un module de collecte collaborative de données textuelles et vocales. L'ensemble des ressources produites, corpus et modèles, a été rendu disponible afin de favoriser la réutilisation scientifique et de contribuer au développement du traitement automatique des langues pour le ruund. À notre connaissance, ce travail constitue le premier système de traduction automatique neuronale et le premier corpus parallèle du ruund rendus publiquement disponibles.

\vspace{7cm}
\textbf{Mots-clés :} Traduction Automatique Neuronale, Langues à Faibles Ressources, Ruund, Corpus Parallèle, Transformer, Fine-tuning, mBART-50, MarianMT, NLLB-200-distilled-600M, Traitement Automatique des Langues.